% ECONOMICS_PROBLEM: {"id": "C73-651", "title": "Karlin's fictitious-play rate conjecture: disproved", "jel_code": "C73", "source_id": "OP-0146"}
% !TeX program = xelatex
\documentclass[11pt,a4paper]{article}
\usepackage[margin=23mm]{geometry}
\usepackage{fontspec}
\setmainfont{Latin Modern Roman}
\usepackage{amsmath,amssymb,mathtools}
\usepackage{xcolor,enumitem,booktabs,longtable,array,xurl,hyperref}
\definecolor{muted}{HTML}{53636C}
\hypersetup{colorlinks=true,urlcolor=blue,linkcolor=blue}
\setlength{\parindent}{0pt}
\setlength{\parskip}{5pt}
\newcommand{\R}{\mathbb R}
\newcommand{\E}{\mathbb E}
\newcommand{\N}{\mathbb N}
\newcommand{\ind}{\mathbf 1}
\newcommand{\Var}{\operatorname{Var}}
\newcommand{\Cov}{\operatorname{Cov}}
\newcommand{\supp}{\operatorname{supp}}
\DeclareMathOperator*{\argmin}{arg\,min}
\DeclareMathOperator*{\argmax}{arg\,max}
\newcommand{\entrymeta}[3]{{\sffamily\small\color{muted}\textbf{\texttt{#1}}\quad|\quad #2\par #3\par}}
\newcommand{\Problemref}[1]{\texttt{#1}}
\newcommand{\JELref}[2]{\texttt{#2} (JEL \texttt{#1})}
\begin{document}
\section*{Karlin's fictitious-play rate conjecture: disproved}
\textbf{Problem ID:} \texttt{C73-651}\quad \textbf{JEL:} \texttt{C73}\par
% BEGIN ECONOMICS STATEMENT
\label{problem:OP-0146}
% GAME_THEORY_PROBLEM: {"local_id": "GT-016", "original_number": 16, "kind": "H", "entry_kind": "statement_only", "published": false, "review_required": true, "category": "game_theory"}
\label{game:GT-016}
{\sffamily\small\color{muted}
\textbf{GT-016} | Learning in zero-sum games\\
\textbf{H} | False historical conjecture; both tie-breaking qualifications must be retained.
\par}
\subsection*{Definitions and mathematical statement}
Let $A\in\mathbb R^{m\times n}$ be the payoff matrix of a maximizing row player against a minimizing column player. Initialize $x^1=e_{i_1}$ and $y^1=e_{j_1}$. Simultaneous fictitious play chooses, for $t\geq1$,
\[
 i_{t+1}\in\operatorname{argmax}_i(Ay^t)_i,\qquad
 j_{t+1}\in\operatorname{argmin}_j((x^t)^{\mathsf T}A)_j,
\]
and updates $x^{t+1}=(tx^t+e_{i_{t+1}})/(t+1)$ and analogously for $y$. Define the duality gap
\[
 g_A(x,y)=\max_i(Ay)_i-\min_j(x^{\mathsf T}A)_j.
\]
The strong historical conjecture asserted that, for every fixed $A$, every initialization, and every legal tie-breaking sequence, there exist finite $C,T_0$ such that $g_A(x^T,y^T)\leq CT^{-1/2}$ for all $T\geq T_0$. This assertion is false. [S1] supplies identity-matrix games and tie-breaking for which the gap decays as slowly as $\Omega(T^{-1/n})$. [S2] supplies a $10\times10$ game with a slower-than-$T^{-1/2}$ rate and no ties after the first step, defeating a weaker tie-breaking qualification as well.
\subsection*{Short English statement}
Ordinary fictitious play need not converge at the conjectured square-root rate, even when later ties are absent.
\subsection*{Variants and scope boundaries}
Some special game classes and specific tie-breaking rules can satisfy the rate. Perturbed best responses, stochastic fictitious play, and anticipatory variants change the update rule. Those variants do not restore the false universal statement for ordinary fictitious play.
\subsection*{Source relationship and status}
The supplied catalog mislabels a false conjecture as open. Its cited perturbation paper [S3] contains a historical convergence-rate remark but does not supersede the explicit counterexamples. This section preserves the original question for audit purposes, without relabelling it as an open problem.
\subsection*{References}
\begingroup\footnotesize\raggedright
{\footnotesize\raggedright
\textbf{[S1]} Constantinos Daskalakis and Qinxuan Pan (2014). \emph{A Counter-Example to Karlin's Strong Conjecture for Fictitious Play}. FOCS 2014; arXiv:1412.4840. Locator: abstract and identity-matrix counterexample. \url{https://arxiv.org/abs/1412.4840}\par
\textbf{[S2]} Yuanhao Wang (2025). \emph{Tie-breaking Agnostic Lower Bound for Fictitious Play}. arXiv:2507.09902. Locator: abstract, Theorem 1 and construction of the augmented game. \url{https://arxiv.org/pdf/2507.09902}\par
\textbf{[S3]} Adam Dziwoki and Rostislav Hor\v{c}\'ik (2025). \emph{Perturbing Best Responses in Zero-Sum Games}. arXiv:2511.12523v1. Locator: ``Related Work'' and Theorem 3; different, perturbed dynamics. \url{https://arxiv.org/html/2511.12523}\par}
\par\endgroup

\subsection*{Common conventions: Source mathematical conventions}

\textbf{Finite games.} For a finite set $A$, $\Delta(A)$ is its probability
simplex. Players choose independent mixed strategies in Nash equilibrium;
correlation is allowed only when explicitly part of the solution concept.
In a game with payoff functions $u_i$ and strategy profile $x$, an additive
$\varepsilon$ Nash equilibrium satisfies
\[
 u_i(x)\geq u_i(x_i',x_{-i})-\varepsilon
 \quad\text{for every player $i$ and every admissible deviation $x_i'$}.
\]
For numerical approximation, payoffs are normalized as locally specified.
An $\varepsilon$ payoff residual is distinct from distance at most
$\varepsilon$ to the set of exact equilibria. Point convergence, convergence
of empirical play, and convergence in distance to an equilibrium set are
also distinct.

\textbf{Computational representations.} Unless stated otherwise, a complexity
question takes rational data with binary encoding; polynomial time is measured
in total bit length. A fixed-accuracy polynomial algorithm is not necessarily
polynomial in $1/\varepsilon$, and neither is necessarily polynomial in
$\log(1/\varepsilon)$. Succinct games and value, demand, or sampling oracles
must declare their interfaces. Exact real solutions require an explicit
representation; an unspecified list of arbitrary real numbers is not a
finite computational output. A claimed algorithmic barrier is meaningful
only for its specified model and reductions.

\textbf{Dynamic games.} Histories, information, observation, and allowable
randomization are part of the model. A stationary strategy depends only on
the current state; a positional strategy in a deterministic graph game
chooses one action at each controlled vertex. Nash equilibrium at an initial
state and equilibrium after every history are different requirements.
An expectation of a pathwise $\liminf$ payoff, a $\liminf$ of expected averages,
a limiting expected average, and a uniform finite-horizon equilibrium are
different criteria. None is silently substituted for another.

\textbf{Allocation and mechanisms.} A complete allocation assigns every item
exactly once unless otherwise stated. Zero-valued goods, negative values,
capacity constraints, feasibility, lotteries, and copies of goods affect
fairness definitions and are specified locally. Dominant-strategy incentive
compatibility, Bayesian incentive compatibility, universal truthfulness,
and truthfulness in expectation impose different inequalities. Whenever
an objective ratio has zero denominator, its convention is stated or those
instances are excluded.

\textbf{Infinite-dimensional models.} Expectations, integrals, stochastic
equations, and optimization objectives require their local regularity,
integrability, filtrations, and admissible domains. A backward--forward
mean-field system is a boundary-value problem; it is not an initial-value
semiflow without an additional construction. Long-time, large-population,
and vanishing-noise limits require an order of limits and a convergence
topology. If a minimum need not be attained, the objective is its infimum
or supremum and the approximation requirement is stated separately.
% END ECONOMICS STATEMENT
\end{document}
